\subsection{LINEAR EQUATIONS OF ORDER n}

One calls a linear differential equation of order n an equation of the form

\[
\mathrm{D} ^ {n} x - a _ {1} (t) \mathrm{D} ^ {n - 1} x - \dots - a _ {n - 1} (t) \mathrm{D} x - a _ {n} (t) x = b (t)\tag{33}
\]

where the \(a_{k}\) ( \(1 \leqslant k \leqslant n\) ) and b are real (complex) functions of the real variable t, defined on an interval J of R. The general procedure of IV, p.~164 shows that this equation is equivalent to the linear system of n first order equations

\[
\left\{ \begin{array}{l l} \frac {d x _ {k}}{d t} = x _ {k + 1} & (1 \leqslant k \leqslant n - 1) \\ \frac {d x _ {n}}{d t} = a _ {1} (t) x _ {n} + a _ {2} (t) x _ {n - 1} + \dots + a _ {n} (t) x _ {1} + b (t) \end{array} \right.\tag{34}
\]

that is to say, to the linear equation

\[
\frac {d \mathbf {x}}{d t} = A (t). \mathbf {x} + \mathbf {b} (t)\tag{35}
\]

where we have put \(\mathbf{x}=(x_{1},x_{2},\ldots,x_{n})\in\mathbf{C}^{n},\mathbf{b}(t)=(0,0,\ldots,0,b(t))\) , and where the matrix \(A(t)\) is defined by

\[
A (t) = \left( \begin{array}{c c c c c} 0 & 1 & 0 & \ldots & 0 \\ 0 & 0 & 1 & \ldots & 0 \\ \vdots & \vdots & \vdots & \ldots & \vdots \\ 0 & 0 & 0 & \ldots & 1 \\ a _ {n} (t) & a _ {n - 1} (t) & a _ {n - 2} (t) & \ldots & a _ {1} (t) \end{array} \right).
\]

The study of the linear equation of order n thus consists of applying the general results above to the particular linear equation (35). For every interval J where the functions \(a_{j}\) ( \(1 \leqslant j \leqslant n\) ) and b are regulated there exists one and only one function u, defined on J, having a continuous derivative of order n - 1, and a regulated derivative of order n on this interval (except at the points of a countable set), satisfying (33) on the complement of a countable subset of J, and such that

\[
u (t _ {0}) = x _ {0}, \qquad \mathrm{D} u (t _ {0}) = x _ {0} ^ {\prime}, \dots , D ^ {n - 1} u (t _ {0}) = x _ {0} ^ {(n - 1)}\tag{36}
\]

where \(t_0\) is an arbitrary point of \(J\) , and \(x_0, x_0', \ldots, x_0^{(n-1)}\) are \(n\) arbitrary complex numbers.

For the \(p\) integrals \(u_{j}\) ( \(1 \leqslant j \leqslant p\) ) of the homogeneous equation

\[
\mathrm{D} ^ {n} x - a _ {1} (t) \mathrm{D} ^ {n - 1} x - \dots - a _ {n - 1} (t) \mathrm{D} x - a _ {n} (t) x = 0\tag{37}
\]

associated with (33) to be linearly independent (in the space \(\mathcal{C}(\mathrm{J};\mathbf{C})\) of continuous maps from J into C, considered as a vector space over C), it is necessary and sufficient that the corresponding p integrals \(\mathbf{u}_{j}=(u_{j},\mathrm{D}u_{j},\ldots,\mathrm{D}^{n-1}u_{j})\) of the homogeneous equation \(d\mathbf{x}/dt=A(t)\mathbf{x}\) should be linearly independent (in the space \(\mathcal{C}(\mathrm{J};\mathbf{C}^{n})\) of continuous maps from J into \(C^{n}\) ). It is clear that this condition is necessary.

Conversely, if there are \(n\) complex constants \(\lambda_j\) , not all zero, such that \(\sum_{j=1}^{n} \lambda_j u_j(t) = 0\) identically on \(J\) , one deduces that \(\sum_{j=1}^{n} \lambda_j D^k u_j(t) = 0\) on \(J\) for every integer \(k\) such that \(1 \leqslant k \leqslant n - 1\) , which means that \(\sum_{j=1}^{n} \lambda_j \mathbf{u}_j(t) = 0\) on \(J\) .

Consequently (IV, p.~180, cor. 1)

PROPOSITION 8. The set of integrals of the homogeneous linear equation (37), defined on J, is a vector space of dimension n over the field C.

Given any n integrals \(u_{j}\) ( \(1 \leqslant j \leqslant n\) ) of equation (37) one calls the Wronskian of this system of integrals the determinant (with respect to the canonical basis of \(C^{n}\) ) of the corresponding n integrals \(u_{j}\) of the equation \(d\mathbf{x}/dt = A(t)\mathbf{x}\) , that is to say, the function

\[
\mathrm{W} (t) = \left| \begin{array}{c c c c} u _ {1} (t) & u _ {2} (t) & \ldots & u _ {n} (t) \\ \mathrm{D} u _ {1} (t) & \mathrm{D} u _ {2} (t) & \ldots & \mathrm{D} u _ {n} (t) \\ \vdots & \vdots & \ldots & \vdots \\ \mathrm{D} ^ {n - 1} u _ {1} (t) & \mathrm{D} ^ {n - 1} u _ {2} (t) & \ldots & \mathrm{D} ^ {n - 1} u _ {n} (t) \end{array} \right|.
\]

For the n integrals \(u_{j}\) to be linearly independent it is necessary and sufficient that \(\mathrm{W}(t) \neq 0\) on J; moreover, it suffices for this that \(\mathrm{W}(t_{0}) \neq 0\) at only one \(t_{0}\) of J (IV, p.~184, prop. 4); further, one has (IV, p.~184, prop. 5)

\[
\mathrm{W} (t) = \mathrm{W} (t _ {0}) \exp \left(\int_ {t _ {0}} ^ {t} a _ {1} (s) d s\right).\tag{38}
\]

We identify the resolvent \(C(t, t_{0})\) of equation (35) with its matrix with respect to the canonical basis of \(C^{n}\) ; the columns \(\mathbf{v}_{j}(t, t_{0})\) ( \(1 \leqslant j \leqslant n\) ) of this matrix are then n linearly independent integrals

\[
\mathbf {v} _ {j} (t, t _ {0}) = (v _ {j} (t, t _ {0}), \mathrm{D} v _ {j} (t, t _ {0}), \dots , \mathrm{D} ^ {n - 1} v _ {j} (t, t _ {0}))
\]

of the homogeneous equation \(d\mathbf{x}/dt = A(t)\) . x which correspond to n linearly independent integrals \(v_{j}(t, t_{0})\) of equation (37) such that

\[
\mathbf {D} ^ {k - 1} v _ {j} (t _ {0}, t _ {0}) = \delta_ {j k}
\]

(Kronecker delta) for \(1 \leqslant j \leqslant n\) , \(1 \leqslant k \leqslant n\) (on agreeing to put \(D^{0}v_{j} = v_{j}\) ). It results in particular that the method of variation of parameters (IV, p.~182) applied to equation (35) here gives as a particular integral of (33), equal to 0 together with its first n - 1 derivatives at the point \(t_{0}\) , the function

\[
w (t) = \int_ {t _ {0}} ^ {t} v _ {n} (t, s) b (s) d s.\tag{39}
\]

In the particular case of the equation \(D^{n}x = b(t)\) the formula (39) again gives the formula expressing the \(n^{th}\) primitive of the regulated function \(b(t)\) which vanishes with its first n - 1 derivatives at the point \(t_{0}\) , namely

\[
w (t) = \int_ {t _ {0}} ^ {t} b (s) \frac {(t - s) ^ {n - 1}}{(n - 1) !} d s
\]

(II, p.~62, formula (19)): the integral of \(D^{n}x = 0\) which vanishes together with its first n - 2 derivatives at the point \(t_{0}\) , and whose \(n - 1^{th}\) derivative is equal to 1 there, is actually the polynomial \((t - t_{0})^{n-1}/(n - 1)!\) .

